\documentclass[10pt]{article}
\usepackage[margin=0.85in]{geometry}
\usepackage{booktabs}
\usepackage{graphicx}
\usepackage{hyperref}
\usepackage{amsmath}

\renewcommand{\arraystretch}{0.9}
\makeatletter
\renewcommand\section{\@startsection{section}{1}{0pt}{1.2ex plus .5ex minus .1ex}{0.6ex plus .1ex}{\normalfont\large\bfseries}}
\renewcommand\subsection{\@startsection{subsection}{2}{0pt}{0.8ex plus .4ex minus .1ex}{0.3ex plus .1ex}{\normalfont\bfseries}}
\makeatother
\setlength{\abovecaptionskip}{2pt}
\setlength{\belowcaptionskip}{2pt}
\setlength{\textfloatsep}{6pt}
\setlength{\intextsep}{4pt}

\title{\large How much can image features be compressed?\\
\normalsize Predicting retrieval loss under random projection from intrinsic dimension}
\author{\small Md Maidul Islam}
\date{\small November 2026}

\begin{document}
\maketitle
\vspace{-3ex}

\section*{Question}

Image features are routinely compressed to save memory and speed up
nearest-neighbour search. This project asks a narrow, measurable version
of a broader question: when a feature vector is compressed by a random
projection to $m$ dimensions, how much nearest-neighbour retrieval
quality is lost, and can that loss be predicted ahead of time from the
feature's \emph{intrinsic dimension} -- how many real degrees of freedom
the data varies along, as opposed to how many numbers it happens to be
stored with? This is a small, exploratory empirical study, not a novel
method and not a state-of-the-art result.

\section*{Methods}

\subsection*{Part A: image and signal processing from scratch}

Convolution/correlation, Gaussian blur, Sobel gradients, Canny edge
detection (non-maximum suppression and double-threshold hysteresis),
histogram equalization, and FFT-domain low/high-pass filtering were all
implemented in NumPy and validated against OpenCV and the convolution
theorem. See \texttt{docs/EXPLAINERS.md} for the formulas and
\texttt{docs/DECISIONS.md} for implementation choices (e.g. that Sobel
uses correlation, not a flipped kernel, matching \texttt{cv2.Sobel}'s
measured behaviour).

\subsection*{Part B: features}

Three feature types were extracted for the fixed 10{,}000-image CIFAR-10
test set (Krizhevsky, 2009): raw pixels (3072-d, globally normalized to
$[0,1]$), HOG (Dalal \& Triggs, 2005; 324-d), and pretrained ResNet-18
(He et al., 2016) penultimate-layer features (512-d, 224px input
resolution, chosen after a speed benchmark -- see Table~\ref{tab:b}).

\subsection*{Part C: random projection and retrieval}

Three random projections were implemented: dense Gaussian, sparse
Achlioptas (2003), and a structured Subsampled Randomized Hadamard
Transform (SRHT), all scaled to be norm-preserving in expectation.
Ground-truth top-10 Euclidean neighbours were computed in each feature's
original space (10{,}000 queries, leave-query-out); after projection to
$m \in \{4,8,16,32,64,128,256\}$ dimensions ($m<d$ only), top-10 was
recomputed and recall@10 measured, averaged over 10 projection seeds with
a 95\% bootstrap confidence interval (1000 resamples) over the pooled
(seed, query) recall values.

Intrinsic dimension was estimated with TwoNN (Facco et al., 2017) and the
Levina \& Bickel (2004) MLE ($k=10$ to $20$, averaged), both validated on
synthetic data of known dimension before use on real features
(Table~\ref{tab:c-id}).

\section*{Results}

\subsection*{Part A validation}

\begin{table}[h]
\centering
\small
\begin{tabular}{lrrr}
\toprule
Check (ours vs OpenCV) & Max abs error & Mean abs error & Pixel agreement \\
\midrule
Canny\_vs\_canny\_edge\_detector & 1.00e+00 & 2.77e-02 & 0.972 \\
GaussianBlur\_vs\_gaussian\_blur & 1.14e-13 & 1.87e-14 & -- \\
Sobel\_vs\_sobel\_gradients & 0.00e+00 & 0.00e+00 & -- \\
equalizeHist\_vs\_histogram\_equalize & 0.00e+00 & 0.00e+00 & 1.000 \\
filter2D\_vs\_correlate2d & 1.71e-13 & 1.54e-14 & -- \\
\bottomrule
\end{tabular}

\caption{Part A: our NumPy implementations vs.\ OpenCV, max over 4 validation images.}
\label{tab:a-validation}
\end{table}

Agreement is at floating-point precision for filtering and Gaussian blur,
exact for Sobel, and 93--99\% pixel agreement for Canny (a different,
also-valid NMS/threshold convention than OpenCV's, not a bug; see
\texttt{docs/DECISIONS.md}). The convolution theorem was confirmed
numerically (max abs error $\sim 10^{-13}$ between spatial and FFT-domain
convolution, kernel sizes 3--31).

\subsection*{Part B features}

\begin{table}[h]
\centering\small
\begin{tabular}{lrr}
\toprule
Feature & Shape & Extraction time (10{,}000 images) \\
\midrule
pixels & (10000, 3072) & 0.0s \\
hog & (10000, 324) & 1.2s \\
resnet & (10000, 512) & 213.4s \\
\bottomrule
\end{tabular}

\hspace{1em}
\begin{tabular}{rr}
\toprule
ResNet-18 input resolution & Estimated time, 10{,}000 images \\
\midrule
32px & 0.9 min \\
64px & 0.9 min \\
224px & 3.6 min \\
\bottomrule
\end{tabular}

\caption{Left: cached feature shapes/extraction time, all 10{,}000 images. Right: ResNet-18 resolution benchmark (500-image subset, extrapolated) -- motivates using 224px, well under the 20-min budget even at the slowest option.}
\label{tab:b}
\end{table}

\subsection*{Part C: intrinsic dimension}

\begin{table}[h]
\centering\small
\begin{tabular}{lrrr}
\toprule
Synthetic case & True dim. & TwoNN estimate & MLE estimate \\
\midrule
5d linear subspace in 100d & 5 & 5.00 & 5.03 \\
10d hypercube in 100d & 10 & 8.75 & 8.38 \\
\bottomrule
\end{tabular}

\hspace{1em}
\begin{tabular}{lrrr}
\toprule
Feature & $d$ & TwoNN $\hat{ID}$ & MLE $\hat{ID}$ \\
\midrule
pixels & 3072 & 31.63 & 25.11 \\
hog & 324 & 39.95 & 36.09 \\
resnet & 512 & 30.09 & 26.77 \\
\bottomrule
\end{tabular}

\caption{Left: ID estimator validation on synthetic data of known dimension (5-d subspace near-exact; 10-d hypercube underestimated by $\sim$1.5--2, a documented boundary-bias effect, not an error). Right: estimated ID by feature type, full set. Per-class estimates (30 combinations) are in \texttt{results/c\_intrinsic\_dim.csv}.}
\label{tab:c-id}
\end{table}

\subsection*{Part C: recall vs.\ projection dimension}

\begin{table}[h]
\centering\small
\begin{tabular}{llrrr}
\toprule
Feature & Projection & recall@10 (m=64) & CI lower & CI upper \\
\midrule
hog & gaussian & 0.196 & 0.195 & 0.196 \\
hog & sparse & 0.195 & 0.194 & 0.195 \\
hog & srht & 0.210 & 0.209 & 0.211 \\
pixels & gaussian & 0.378 & 0.377 & 0.380 \\
pixels & sparse & 0.375 & 0.374 & 0.376 \\
pixels & srht & 0.383 & 0.381 & 0.384 \\
resnet & gaussian & 0.297 & 0.296 & 0.298 \\
resnet & sparse & 0.300 & 0.299 & 0.301 \\
resnet & srht & 0.322 & 0.321 & 0.323 \\
\bottomrule
\end{tabular}

\caption{recall@10 at $m=64$, all (feature, projection) pairs, mean with 95\% bootstrap CI. Full curves: \texttt{results/figures/c\_recall\_vs\_m\_*.png}. All three projections perform similarly for pixels/HOG; for ResNet, SRHT noticeably outperforms Gaussian/sparse at larger $m$ ($\approx$0.69 vs.\ $\approx$0.58 at $m=256$), reported but not fully explained.}
\label{tab:c-recall-m64}
\end{table}

\subsection*{Part C: collapse hypothesis}

\begin{table}[h]
\centering\small
\begin{tabular}{lr}
\toprule
Comparison & Value \\
\midrule
m & 0.0425 \\
m\_over\_d & 0.1897 \\
m\_over\_id & 0.0305 \\
spearman\_id\_vs\_recall\_loss\_at\_m=64\_n=3 & 0.5000 \\
\bottomrule
\end{tabular}

\caption{Collapse-spread statistic by $x$-axis (lower = more collapse), plus Spearman correlation. $m/\hat{ID}$ gives a slightly lower spread than raw $m$ (0.031 vs.\ 0.043), but the three feature types' TwoNN estimates are themselves fairly close (31.6, 40.0, 30.1), leaving limited room for $m/\hat{ID}$ to differ from $m$; this is weak, preliminary evidence at best, based on only three curves with no uncertainty estimate on the statistic itself. The $n=3$ Spearman check is reported for completeness, not as evidence.}
\label{tab:c-collapse}
\end{table}

\section*{Limitations}

\begin{itemize}
\setlength{\itemsep}{1pt}
\setlength{\parsep}{0pt}
\setlength{\topsep}{2pt}
\item Single dataset (CIFAR-10), a fixed 10{,}000-image subset, one
  random seed for data splits.
\item CPU-only, which bounded some choices (e.g.\ ResNet-18 instead of a
  larger backbone) but was not otherwise a binding constraint -- every
  experiment here, including the slowest (18.5 minutes), finished well
  within the hour.
\item The collapse analysis uses only feature-type-level recall curves,
  not separate recall sweeps per class, because the per-class version
  would have cost roughly 10$\times$ Part C's measured runtime; per-class
  \emph{intrinsic dimension} estimates were still computed in full
  (\texttt{results/c\_intrinsic\_dim.csv}).
\item The hypothesis (recall vs.\ $m/\hat{ID}$ collapses across feature
  types) is exploratory and received only weak, preliminary support, not
  confirmation: the three feature types' intrinsic dimension estimates
  are themselves fairly close together, leaving limited room for
  $m/\hat{ID}$ to differ meaningfully from raw $m$. The Spearman
  correlation check is underpowered ($n=3$).
\item Canny does not pixel-match OpenCV exactly, by design (a different,
  also-standard NMS/threshold convention); this is a known, measured,
  and documented difference, not an unverified assumption.
\end{itemize}

\section*{Reproducibility}

Every number in this report is generated from a CSV in \texttt{results/}
by \texttt{report/make\_tables.py}; nothing above is typed by hand.
Seed 42 is used everywhere a seed applies. Machine/package versions are
recorded in \texttt{results/environment.txt}. To reproduce from a clean
clone:

\begin{verbatim}
./run_all.sh
\end{verbatim}

This recreates the virtual environment, re-downloads CIFAR-10 if needed,
re-runs every experiment in order, and regenerates every CSV, figure, and
this report's tables. See \texttt{README.md} for the exact commands to
run individual stages and \texttt{docs/DECISIONS.md} for every parameter
choice and the reasoning behind it.

\end{document}
